% GENERATED FILE. Edit canonical lessons, then run npm run generate:books.
% canonical-source-hash: fnv1a64:00b44f28369fda23
% canonical-lessons: GE-C419-zwischenlandung, GE-C419-zeitverschiebung, GE-C419-verspaeten, GE-C419-festsitzen, GE-C419-notlanden

\chapter{Warten am Flughafen}
\label{ch:ge-warten-am-flughafen}
\hlchaptermodality{419}

\begin{chapteropening}
\textbf{By the end of this chapter:} \emph{I can talk about delays: the stopover, the time difference, running late, being stuck and an emergency landing.}

Complete the last lesson of chapter 419: I can talk about delays: the stopover, the time difference, running late, being stuck and an emergency landing.
\end{chapteropening}

\section[die Zwischenlandung]{die Zwischenlandung — the stopover}

\label{lesson:GE-C419-zwischenlandung}



\begin{quote}
Before the new one: say the German for to take off, then the German for the landing.
\end{quote}

\subsection*{You'll want to know: die Zwischenlandung}
\textbf{die Zwischenlandung} — \textquotedblleft{}the stopover\textquotedblright{}.

\begin{quote}
Der Flug hat eine Zwischenlandung von drei Stunden.
\end{quote}

\begin{quote}
\emph{The flight has a three-hour stopover.}
\end{quote}

Plural: \emph{die Zwischenlandungen}.

\begin{cousinweb}
\emph{Zwischenlandung} is built from \emph{zwischen} 'between', inherited from Old High German \emph{in zwiskēn} 'between two' (EtymDB), and \emph{Landung} 'landing', from \emph{Land}.
\end{cousinweb}

\begin{grammarlens}[title={What you've built}]
One more word for telling, explaining and describing at B1.
\end{grammarlens}

\subsection*{Your turn}
\begin{itemize}
\raggedright
  \item \emph{Say it:} \emph{die Zwischenlandung}
  \item \emph{Say it:} \emph{die Zwischenlandung}, once more
  \item \emph{Say it:} a sentence of your own with \emph{die Zwischenlandung}
  \item \emph{Recall:} say the German for to take off, then the German for the landing, then say \emph{die Zwischenlandung} again
\end{itemize}

\begin{tcolorbox}[breakable,colback=teal!4,colframe=teal!35!black,title={Before you move on}]
What does die Zwischenlandung mean? (\textquotedblleft{}The stopover\textquotedblright{}.) Say it once more.
\end{tcolorbox}
\section[die Zeitverschiebung]{die Zeitverschiebung — the time difference; jet lag}

\label{lesson:GE-C419-zeitverschiebung}



\begin{quote}
Before the new one: say the German for the landing, then the German for the stopover.
\end{quote}

\subsection*{You'll want to know: die Zeitverschiebung}
\textbf{die Zeitverschiebung} — \textquotedblleft{}the time difference; jet lag\textquotedblright{}.

\begin{quote}
Wegen der Zeitverschiebung ist er um vier Uhr morgens wach.
\end{quote}

\begin{quote}
\emph{Because of the time difference he is awake at four in the morning.}
\end{quote}

Plural: \emph{die Zeitverschiebungen}. People also say \emph{der Jetlag}.

\begin{cousinweb}
\emph{Zeitverschiebung} is built from \emph{Zeit} 'time', inherited from Proto-Germanic \emph{tīdiz} (EtymDB), and \emph{Verschiebung} 'shift', from \emph{schieben} 'to push', inherited from Old High German \emph{skioban} (EtymDB).
\end{cousinweb}

\begin{grammarlens}[title={What you've built}]
One more word for telling, explaining and describing at B1.
\end{grammarlens}

\subsection*{Your turn}
\begin{itemize}
\raggedright
  \item \emph{Say it:} \emph{die Zeitverschiebung}
  \item \emph{Say it:} \emph{die Zeitverschiebung}, once more
  \item \emph{Say it:} a sentence of your own with \emph{die Zeitverschiebung}
  \item \emph{Recall:} say the German for the landing, then the German for the stopover, then say \emph{die Zeitverschiebung} again
\end{itemize}

\begin{tcolorbox}[breakable,colback=teal!4,colframe=teal!35!black,title={Before you move on}]
What does die Zeitverschiebung mean? (\textquotedblleft{}The time difference; jet lag\textquotedblright{}.) Say it once more.
\end{tcolorbox}
\section[verspäten]{verspäten — to be late, to run late (mostly reflexive: sich verspäten)}

\label{lesson:GE-C419-verspaeten}



\begin{quote}
Before the new one: say the German for the stopover, then the German for the time difference.
\end{quote}

\subsection*{You'll want to know: verspäten}
\textbf{verspäten} — \textquotedblleft{}to be late, to run late (mostly reflexive: sich verspäten)\textquotedblright{}.

\begin{quote}
Der Flug verspätet sich um zwei Stunden.
\end{quote}

\begin{quote}
\emph{The flight is two hours late.}
\end{quote}

Mostly reflexive: \emph{ich verspäte mich}. \emph{Die Verspätung} — the delay.

\begin{cousinweb}
\emph{verspäten} is built from \emph{ver-} and \emph{spät} 'late', an inherited West Germanic word (EtymDB lists \emph{spät}).
\end{cousinweb}

\begin{grammarlens}[title={What you've built}]
One more word for telling, explaining and describing at B1.
\end{grammarlens}

\subsection*{Your turn}
\begin{itemize}
\raggedright
  \item \emph{Say it:} \emph{verspäten}
  \item \emph{Say it:} \emph{verspäten}, once more
  \item \emph{Say it:} a sentence of your own with \emph{verspäten}
  \item \emph{Recall:} say the German for the stopover, then the German for the time difference, then say \emph{verspäten} again
\end{itemize}

\begin{tcolorbox}[breakable,colback=teal!4,colframe=teal!35!black,title={Before you move on}]
What does verspäten mean? (\textquotedblleft{}To be late, to run late (mostly reflexive: sich verspäten)\textquotedblright{}.) Say it once more.
\end{tcolorbox}
\section[festsitzen]{festsitzen — to be stuck}

\label{lesson:GE-C419-festsitzen}



\begin{quote}
Before the new one: say the German for the time difference, then the German for to be late.
\end{quote}

\subsection*{You'll want to know: festsitzen}
\textbf{festsitzen} — \textquotedblleft{}to be stuck\textquotedblright{}.

\begin{quote}
Wegen des Schnees sitzen wir am Flughafen fest.
\end{quote}

\begin{quote}
\emph{Because of the snow we are stuck at the airport.}
\end{quote}

Separable: \emph{wir sitzen fest}. Perfekt with \emph{haben}, in the south also \emph{sein}: \emph{wir haben festgesessen}.

\begin{cousinweb}
\emph{festsitzen} is built from \emph{fest} 'firm, fixed' and \emph{sitzen} 'to sit', inherited from Proto-Germanic \emph{sitjaną} (EtymDB).
\end{cousinweb}

\begin{grammarlens}[title={What you've built}]
One more word for telling, explaining and describing at B1.
\end{grammarlens}

\subsection*{Your turn}
\begin{itemize}
\raggedright
  \item \emph{Say it:} \emph{festsitzen}
  \item \emph{Say it:} \emph{festsitzen}, once more
  \item \emph{Say it:} a sentence of your own with \emph{festsitzen}
  \item \emph{Recall:} say the German for the time difference, then the German for to be late, then say \emph{festsitzen} again
\end{itemize}

\begin{tcolorbox}[breakable,colback=teal!4,colframe=teal!35!black,title={Before you move on}]
What does festsitzen mean? (\textquotedblleft{}To be stuck\textquotedblright{}.) Say it once more.
\end{tcolorbox}
\section[notlanden]{notlanden — to make an emergency landing}

\label{lesson:GE-C419-notlanden}



\begin{quote}
Before the new one: say the German for to be late, then the German for to be stuck.
\end{quote}

\subsection*{You'll want to know: notlanden}
\textbf{notlanden} — \textquotedblleft{}to make an emergency landing\textquotedblright{}.

\begin{quote}
Das kleine Flugzeug musste auf einem Feld notlanden.
\end{quote}

\begin{quote}
\emph{The small plane had to make an emergency landing in a field.}
\end{quote}

Perfekt with \emph{sein}: \emph{das Flugzeug ist notgelandet}. \emph{Die Notlandung} — the emergency landing.

\begin{cousinweb}
\emph{notlanden} is built from \emph{Not} 'need, distress', an inherited word that is a cousin of English \emph{need} (EtymDB), and \emph{landen} 'to land', from \emph{Land}.
\end{cousinweb}

\begin{grammarlens}[title={What you've built}]
One more word for telling, explaining and describing at B1.
\end{grammarlens}

\subsection*{Your turn}
\begin{itemize}
\raggedright
  \item \emph{Say it:} \emph{notlanden}
  \item \emph{Say it:} \emph{notlanden}, once more
  \item \emph{Say it:} a sentence of your own with \emph{notlanden}
  \item \emph{Recall:} say the German for to be late, then the German for to be stuck, then say \emph{notlanden} again
\end{itemize}

\begin{tcolorbox}[breakable,colback=teal!4,colframe=teal!35!black,title={Before you move on}]
What does notlanden mean? (\textquotedblleft{}To make an emergency landing\textquotedblright{}.) Say it once more.
\end{tcolorbox}
